% eksperymentalne tłumaczenie na włoski

%

\subsection{Aksjomaty przystawania kątów} % lang-pl
\index{aksjomat!przystawania kątów}% % lang-pl
\subsection{Assiomi di congruenza degli angoli} % lang-it
\index{assioma!di congruenza degli angoli}% % lang-it

Dokładnie jak dla odcinków, wprowadzamy niezdefiniowaną relację przystawania kątów, oznaczaną znowu przez $\cong$, ponieważ nie prowadzi to do nieporozumień. % lang-pl
Esattamente come per i segmenti, introduciamo una relazione indefinita di congruenza tra angoli, indicata ancora con $\cong$, poiché ciò non genera ambiguità. % lang-it

\begin{axiom}
[przystawania, C4] % lang-pl
[di congruenza, C4] % lang-it
    Dla każdego kąta $\angle BAC$ i każdej $DF$ półprostej istnieje dokładnie jedna półprosta $DE$ po ustalonej stronie prostej $DF$ taka, że $\angle BAC \cong \angle EDF$. % lang-pl
    Per ogni angolo $\angle BAC$ e per ogni semiretta $DF$ esiste esattamente una semiretta $DE$ situata da un lato prefissato della retta $DF$ tale che $\angle BAC \cong \angle EDF$. % lang-it
\end{axiom}

Możemy traktować ten aksjomat jako odpowiednik cyrkla: pozwala przenosić kąty tak, jak Euklides (I.23). % lang-pl
\index{cyrkiel} % lang-pl
Possiamo considerare questo assioma come l'equivalente del compasso: permette di trasportare gli angoli come fa Euclide in (I.23). % lang-it
\index{compasso} % lang-it

\begin{axiom}
[przystawania, C5] % lang-pl
[di congruenza, C5] % lang-it
    Niech $\alpha, \beta, \gamma$ będą kątami. % lang-pl
    Jeśli $\alpha \cong \beta$ oraz $\alpha \cong \gamma$, to $\beta \cong \gamma$. % lang-pl
    Każdy kąt przystaje do siebie: $\alpha \cong \alpha$. % lang-pl
    Siano $\alpha$, $\beta$, $\gamma$ angoli. % lang-it
    Se $\alpha \cong \beta$ e $\alpha \cong \gamma$, allora $\beta \cong \gamma$. % lang-it
    Ogni angolo è congruente a sé stesso: $\alpha \cong \alpha$. % lang-it
\end{axiom}

Mamy wreszcie cechę przystawania bok-kąt-bok, z której można wyprowadzić dodawanie kątów. % lang-pl
\index{cecha przystawania!bok-kąt-bok}% % lang-pl
Inne cechy przystawania to bok-bok-bok i kąt-bok-kąt. % lang-pl
Piszą o nich Hartshorne \cite[s. 99]{hartshorne2000}, Neugebauer \cite[s. 8, 11]{neugebauer_2018}, Pompe \cite[s. 25]{pompe_2022}. % lang-pl
Disponiamo infine del criterio di congruenza lato-angolo-lato, dal quale si può derivare l'addizione degli angoli. % lang-it
\index{criterio di congruenza!lato-angolo-lato}% % lang-it
Altri criteri di congruenza sono lato-lato-lato e angolo-lato-angolo. % lang-it
Ne scrivono Hartshorne \cite[p. 99]{hartshorne2000}, Neugebauer \cite[p. 8, 11]{neugebauer_2018} e Pompe \cite[p. 25]{pompe_2022}. % lang-it

\begin{axiom}
[przystawania, C6] % lang-pl
[di congruenza, C6] % lang-it
    Niech $ABC$ i $DEF$ będą dwoma trójkątami takimi, że $AB \cong DE$, $AC \cong DF$ i $\angle BAC \cong \angle EDF$. % lang-pl
    Wtedy pozostałe boki i kąty są przystające, a razem z nimi całe trójkąty: $BC \cong EF$, $\angle ABC \cong \angle DEF$ i $\angle ACB \cong \angle DFE$. % lang-pl
    Siano $ABC$ e $DEF$ due triangoli tali che $AB \cong DE$, $AC \cong DF$ e $\angle BAC \cong \angle EDF$. % lang-it
    Allora anche i restanti lati e angoli sono congruenti, e con essi i triangoli interi: $BC \cong EF$, $\angle ABC \cong \angle DEF$ e $\angle ACB \cong \angle DFE$. % lang-it
\end{axiom}

Euklides uda, że dowodzi aksjomatu C6 metodą ,,nakładania'' figur, ale Hilbert znajdzie model będący świadkiem, że nie wynika on z poprzednich aksjomatów; patrz \cite[paragraf 11]{hilbert_1988}, jak wiemy od Greenberga \cite[s. 200]{greenberg_2010}. % lang-pl
(W obronie Euklidesa należy dodać, że metoda nakładania pojawia się dwa razy: w (I.4) oraz (I.8); gdyby była słuszna, postulat 4, że wszystkie kąty proste są równe, nie byłby potrzebny!). % lang-pl
\index{kąt!prosty}% % lang-pl
Patrz też do Hartshorne'a \cite[s. 148-158]{hartshorne2000}. % lang-pl
Euclide fingerà di dimostrare l'assioma C6 mediante il metodo della ``sovrapposizione'' delle figure, ma Hilbert costruirà un modello che mostra come esso non segua dagli assiomi precedenti; si veda \cite[paragrafo 11]{hilbert_1988}, come sappiamo da Greenberg \cite[p. 200]{greenberg_2010}. % lang-it
(In difesa di Euclide va aggiunto che il metodo della sovrapposizione compare due volte, in (I.4) e (I.8); se fosse valido, il postulato 4, secondo cui tutti gli angoli retti sono uguali, non sarebbe necessario!). % lang-it
\index{angolo!retto}% % lang-it
Si veda anche Hartshorne \cite[p. 148-158]{hartshorne2000}. % lang-it

Jeżeli chodzi o dodawanie kątów, musimy być ostrożni. % lang-pl
Suma dwóch kątów może okazać się prostą (lub ,,dwoma kątami prostymi'' jak pisze Euklides) albo mieć miarę większą niż $\pi$, wtedy składniki sumy nie znajdują się we wnętrzu otrzymanego kąta. % lang-pl
Kąty można także odejmować. % lang-pl
Per quanto riguarda l'addizione degli angoli, occorre essere prudenti. % lang-it
La somma di due angoli può risultare un angolo piatto (o ``due angoli retti'', come scrive Euclide) oppure avere ampiezza maggiore di $\pi$; in tal caso gli addendi non si trovano all'interno dell'angolo ottenuto. % lang-it
Gli angoli possono anche essere sottratti. % lang-it

\begin{definition}
[kąt przyległy] % lang-pl
[angolo adiacente] % lang-it
    Niech $\angle BAC$ będzie kątem, zaś $D$ punktem na prostej $AC$ po przeciwnej stronie $A$ niż $C$. % lang-pl
    Wtedy kąty $\angle BAC$ i $\angle BAD$ nazywamy przyległymi. % lang-pl
    \index{kąt!przyległy} % lang-pl
    Sia $\angle BAC$ un angolo e sia $D$ un punto sulla retta $AC$, dal lato opposto di $A$ rispetto a $C$. % lang-it
    Allora gli angoli $\angle BAC$ e $\angle BAD$ si dicono adiacenti. % lang-it
    \index{angolo!adiacente} % lang-it
\end{definition}

Warto teraz spojrzeć na (I.13) Euklidesa. % lang-pl
(Są jeszcze kąty naprzemianległe i odpowiadające...) % lang-pl
Vale ora la pena di osservare la proposizione (I.13) di Euclide. % lang-it
(Esistono anche gli angoli alterni interni e quelli corrispondenti...) % lang-it

\begin{figure}[H] \centering
...
    \caption{kąty przyległe ($\alpha$ i $\beta$) oraz wierzchołkowe ($\alpha$ i $\gamma$)} % lang-pl
    \caption{angoli adiacenti ($\alpha$ e $\beta$) e opposti al vertice ($\alpha$ e $\gamma$)} % lang-it
\end{figure}

\begin{proposition}
\index{kąt!wierzchołkowy}% % lang-pl
    Kąty wierzchołkowe (para kątów przy dwóch przecinających się prostych, które nie są przyległe) są przystające. % lang-pl
\index{angolo!opposto al vertice}% % lang-it
    Gli angoli opposti al vertice (la coppia di angoli formata da due rette incidenti che non sono adiacenti) sono congruenti. % lang-it
\end{proposition}

Piszą o tym Bogdańska, Neugebauer \cite[s. 6]{neugebauer_2018}. % lang-pl
Pompe \cite[s. 5]{pompe_2022} pisze, że kąty naprzemianległe są równej miary wtedy i tylko wtedy, gdy dwie stosowne proste są równoległe. % lang-pl
Ne scrivono Bogdańska e Neugebauer \cite[p. 6]{neugebauer_2018}. % lang-it
Pompe \cite[p. 5]{pompe_2022} scrive che gli angoli alterni interni hanno la stessa ampiezza se e solo se le due rette considerate sono parallele. % lang-it

% odpowiadające i naprzemianległe.
% https://en.wikipedia.org/wiki/Saccheri-Legendre_theorem

Istnieje odpowiednik relacji między odcinkami ,,bycia mniejszym'' dla kątów: % lang-pl
Esiste anche per gli angoli un analogo della relazione ``essere minore'' definita per i segmenti: % lang-it

\begin{definition}
    Niech $\angle BAC$ i $\angle EDF$ będą kątami. % lang-pl
    Jeśli istnieje półprosta $DG$ wewnątrz kąta $EDF$ taka, że $\angle BAC \cong \angle GDF$, to powiemy, że kąt $\angle BAC$ jest mniejszy niż kąt $\angle EDF$ (a ten drugi jest większy niż pierwszy). % lang-pl
    Siano $\angle BAC$ e $\angle EDF$ due angoli. % lang-it
    Se esiste una semiretta $DG$ interna all'angolo $EDF$ tale che $\angle BAC \cong \angle GDF$, diremo che l'angolo $\angle BAC$ è minore dell'angolo $\angle EDF$ (e quest'ultimo maggiore del primo). % lang-it
\end{definition}

\begin{definition}
    Kąt, który przystaje do kąta przyległego, nazywamy prostym. % lang-pl
    Dwie proste, które przecinają się i wyznaczają tam cztery kąty proste, będziemy nazywać prostopadłymi. % lang-pl
    Un angolo congruente al proprio angolo adiacente si chiama retto. % lang-it
    Due rette che si intersecano formando quattro angoli retti si dicono perpendicolari. % lang-it
\end{definition}

Nie jest ważne, który kąt przyległy wybierzemy, ponieważ kąty wierzchołkowe są przystające. % lang-pl
Non importa quale angolo adiacente si scelga, poiché gli angoli opposti al vertice sono congruenti. % lang-it

%